\section{Introduction}
\label{sec:introduction}

Urbanization and rapid motorization have placed severe pressure on road transportation networks worldwide, leading to chronic traffic congestion, safety hazards, and environmental degradation in developing nations. To mitigate these issues, municipal authorities are deploying Intelligent Transportation Systems (ITS) that leverage real-time traffic monitoring to optimize signal timing, coordinate emergency responses, and measure vehicular flow by category~\cite{Zhou2025TrafficSurveillanceSurvey}. In Peru, the Ministry of Transportation and Communications (MTC) has identified automated vehicle detection and classification as a priority tool for modernizing urban mobility management, formalizing this need through the SMART Challenge 2026~\cite{MTC2026SmartChallenge}. The challenge requires detecting and classifying nine vehicle types (including rare categories such as articulated buses and mototaxis) in every frame of intersection footage, using oriented bounding boxes (OBB) and evaluating with a class-balanced metric that treats each category equally regardless of its frequency.

The widespread availability of closed-circuit television (CCTV) cameras has shifted traffic surveillance toward computer-vision-based methods, which offer rich spatial descriptions of intersection behavior without the infrastructure cost of physical induction loops~\cite{Azfar2024ComplexTrafficReview}. Deep learning techniques, particularly Convolutional Neural Networks (CNNs), have replaced handcrafted feature pipelines by learning hierarchical visual representations directly from image data~\cite{Zou2023ObjectDetectionSurvey}. Generalizing these models across diverse camera angles, illumination conditions, and vehicle populations (especially when training and deployment domains differ) remains an open challenge~\cite{Koga2021VehicleDomainAdaptation}. Multi-view deployments, where several cameras observe the same intersection from different angles, further demand that detectors be robust to viewpoint changes and partial occlusion~\cite{Vora2023GeneralizationMultiview}.

Standard object detectors represent vehicle positions using horizontal bounding boxes (HBB) (rectangles aligned with the image axes). In urban intersections, where vehicles travel diagonally and turn through arbitrary angles, HBB representations enclose substantial background pixels and create artificial overlaps between neighboring boxes~\cite{Wang2025OBBSurvey}. Oriented bounding boxes (OBB), parameterized by center, dimensions, and a rotation angle, produce a geometrically tighter fit that is critical for downstream metrics sensitive to localization precision, such as Time-to-Collision (TTC)~\cite{Ahmed2025}. Foundational approaches to deep OBB detection include region proposal networks modified for rotated candidate boxes~\cite{Xie2021OrientedRCNN,Ding2022DOTA} and rotation-equivariant feature extractors that handle arbitrary orientations within the network backbone~\cite{Han2021ReDet}.

Real-time OBB detection at urban intersections must operate under strict computational constraints. Because streaming high-definition video to centralized cloud infrastructure is prohibitively expensive for most municipalities, inference must run on edge-computing hardware deployed locally~\cite{Balamuralidhar2021}. Single-stage detectors from the YOLO (You Only Look Once) family achieve the best-known accuracy-latency trade-off under these constraints~\cite{Wang2024YOLOv10}. The latest generation, YOLO26, further improves on this trade-off by eliminating Non-Maximum Suppression (NMS) through a dual-head architecture and introducing STAL label assignment, which guarantees positive anchor coverage for small objects~\cite{Jocher2026}. Task-wise sampling convolutions have also been proposed to align classification and localization features at low computational cost~\cite{Huang2025TSConv}.

Peruvian urban traffic is characterized by a highly heterogeneous vehicle fleet where cars dominate the count distribution and compact vehicles (motorcycles, mototaxis, combis) are underrepresented both on the road and in training data. This \textit{class imbalance} causes gradient updates during training to be dominated by the majority class, leading neural networks to learn weak representations for rare categories~\cite{Espinosa2020}. Compact and distant vehicles are additionally vulnerable to feature loss in deep downsampling stacks, as their small spatial footprint occupies too few feature-map cells to survive repeated pooling operations~\cite{Cheng2023SmallObjectSurvey}. In oriented detection specifically, small angular errors compound the problem: a slight orientation mismatch between a predicted OBB and its ground-truth counterpart produces a rotated IoU (rIoU) far lower than the equivalent error would cause in axis-aligned detection~\cite{Zand2022OBB}.

A critical and previously undocumented limitation of the SMART Challenge 2026 dataset is that only \textit{moving} vehicles are annotated: stationary and parked vehicles (fully visible and unoccluded in the frame) receive no bounding box. This annotation convention reflects the MTC's intent to measure active traffic flow, but it creates a fundamental mismatch for any OBB detector, which learns to identify vehicles by visual appearance and has no mechanism to infer motion state from a single frame. A well-trained detector will correctly localize a parked car, but the ground truth will score it as a false positive, uniformly degrading per-class Average Precision across all rIoU thresholds for every architecture evaluated. No training strategy, no architecture choice, and no amount of additional data can resolve this artefact without incorporating temporal information: the motion state of a vehicle can only be determined by comparing its centroid position across consecutive frames within the same clip~\cite{FernandezRodriguez2023,Chen2025}. This problem is not addressed in any of the related work reviewed in this study.

Together, these challenges reveal a gap that no single prior study covers simultaneously: (i) no existing published work evaluates YOLO26-OBB on a real-world urban vehicle dataset, nor compares it fairly against its immediate predecessor YOLO11-OBB; (ii) none of the reviewed works addresses the combination of OBB detection with explicit mitigation of class imbalance for rare vehicle types in a Latin American urban context; and (iii) no prior study documents or corrects for the annotation bias introduced by motion-only labeling, nor quantifies how much of the measured AP gap is attributable to this bias rather than detector performance.

This work addresses the identified gap through the following contributions:
\begin{enumerate}
  \item A \textbf{temporal motion filter} that uses inter-frame centroid displacement within a short clip window to suppress false positives from stationary vehicles, enabling fair comparison of OBB detectors under the motion-only annotation convention.
  \item An \textbf{ablation study} comparing YOLO11-OBB and YOLO26-OBB with and without class-weighted focal loss, isolating the contribution of each component (architecture, angular parameterization, label assignment, and loss function) to the official Macro AP-rIoU\allowbreak@[0.50:0.80] metric.
  \item A \textbf{publicly available implementation} of Macro AP-rIoU\allowbreak@[0.50:0.80] using greedy one-to-one per-class matching over rotated IoU thresholds from 0.50 to 0.80, which does not currently exist as a validated open-source tool.
  \item An \textbf{edge-hardware benchmark} reporting the precision\allowbreak-vs\allowbreak-speed curve of the best model under CPU-only constraints, providing concrete evidence of deployment viability for Peruvian municipalities.
\end{enumerate}

The remainder of this paper is structured as follows. Section~\ref{sec:related_work} reviews related work in oriented vehicle detection and tracking. Section~\ref{sec:material_methods} describes the dataset, the temporal motion filter, and the experimental setup. Section~\ref{sec:results} presents quantitative results. Section~\ref{sec:discussion} discusses findings and limitations, and Section~\ref{sec:conclusions} concludes the study.
